% Shared typography for the role-specific, one-page application résumés.
\documentclass[letterpaper,10pt]{article}
\usepackage{fontspec}
\setmainfont{texgyrepagella}[Extension=.otf,UprightFont=*-regular,BoldFont=*-bold,ItalicFont=*-italic,BoldItalicFont=*-bolditalic,Ligatures=TeX,RawFeature={-liga,-clig}]
\usepackage[margin=0.55in]{geometry}
\usepackage{enumitem}
\usepackage{xcolor}
\usepackage{soul}
\usepackage{hyperref}
\usepackage{microtype}
\usepackage{titlesec}
\definecolor{linkblue}{HTML}{3657D6}
\hypersetup{colorlinks=true,urlcolor=linkblue,linkcolor=linkblue,pdfauthor={Jeet Shah}}
\setul{0.18ex}{0.45pt}
\setulcolor{linkblue}
\let\originalhref\href
\renewcommand{\href}[2]{\originalhref{#1}{\textcolor{linkblue}{\ul{#2}}}}
\pagestyle{empty}
\setlength{\parindent}{0pt}
\setlength{\parskip}{0pt}
\emergencystretch=1em
\raggedright
\setlist[itemize]{leftmargin=1.2em,itemsep=1pt,topsep=2pt,parsep=0pt,partopsep=0pt}
\titleformat{\section}{\normalsize\bfseries}{}{0pt}{}[\vspace{1pt}\titlerule]
\titlespacing{\section}{0pt}{6pt}{3pt}
\newcommand{\entry}[3]{\textbf{#1}\hfill\textit{#2}\par #3\par}
\newcommand{\project}[2]{\textbf{#1}\hfill #2\par}
% Plain-text metadata is also read by the website's resume catalog.
\newcommand{\resumemetadata}[3]{%
  \def\resumetitle{#1}%
  \def\resumeupdated{#2}%
  \hypersetup{pdftitle={Jeet Shah - #1 Resume},pdfsubject={#3 Updated #2}}%
}
\newcommand{\resumeheader}{%
  \begin{center}
    {\fontsize{23}{25}\selectfont\bfseries Jeet Shah}\\[3pt]
    \resumetitle\\[4pt]
    \href{mailto:jeetsh4h@gmail.com}{jeetsh4h@gmail.com} \enspace|\enspace
    \href{https://jeetsh4h.dev}{jeetsh4h.dev} \enspace|\enspace
    \href{https://github.com/jeetsh4h}{github.com/jeetsh4h} \enspace|\enspace
    \href{https://www.linkedin.com/in/jeetsh4h}{linkedin.com/in/jeetsh4h}
  \end{center}
}

\resumemetadata{Technical Researcher}{2026-09-30}{Research software, precipitation nowcasting, HPC work, and preprints.}
\titlespacing{\section}{0pt}{8pt}{3pt}
\begin{document}
\resumeheader
\section{Education}
\entry{Columbia University}{Sep 2026 -- Expected graduation: 12/27}{M.S. Computer Science, Network Systems Pathway · New York, NY}
Current coursework: Networks for AI, Analysis of Algorithms.\par
\vspace{3pt}
\entry{FLAME University}{Aug 2021 -- May 2025}{B.Sc. (Hons.) Computer Science; PG Diploma in Interdisciplinary Studies · Pune, India}
CGPA: 8.9/10; merit scholarships covering 25\% (2021--2024) and 60\% (2024--2025) of tuition.\\
Selected completed coursework: Advanced Machine Learning, Applied Multivariate Statistics, Numerical Methods.
\section{Research Experience}
\entry{Centre for Interdisciplinary AI, FLAME University}{Jan 2024 -- May 2025}{Research Assistant}
\begin{itemize}
  \item Co-developed ConvGRU, conditional-GAN, and ConvLSTM training pipelines and satellite-data ingestion tooling for precipitation nowcasting.
  \item Administered a multi-user GPU server: accounts, drivers, storage, and isolated training environments.
\end{itemize}
\entry{Space Applications Centre, ISRO}{May 2023 -- Aug 2023}{Research Intern}
\begin{itemize}
  \item Developed ConvLSTM precipitation-nowcasting models with INSAT-3D satellite data in an air-gapped high-performance computing environment.
\end{itemize}
\section{Selected Research Outputs}
\textbf{Weather4Cast, team kaubega:} First on the 2024 core-challenge leaderboard; second in the 2025 cumulative-rainfall task.\par
\project{INSAT-3D Precipitation Nowcasting}{\href{https://github.com/jeetsh4h/DISS384}{Thesis \& code}}
\begin{itemize}
  \item Built satellite-data caching, temporal-window generation, and ConvLSTM training/evaluation pipelines in Python, TensorFlow/Keras, OpenCV, and HDF5.
  \item Reduced six-hour forecast RMSE by 30.7\% against a Lucas--Kanade optical-flow baseline; undergraduate thesis supervised by Kaushik Gopalan (2025).
\end{itemize}
\textbf{Preprints and Presentations}\par
\begin{itemize}
  \item Co-author, \href{https://arxiv.org/abs/2511.11197}{Computationally-efficient deep learning models for nowcasting of precipitation} (Weather4Cast 2025), arXiv:2511.11197. \textit{Preprint.}
  \item Co-author, \href{https://arxiv.org/abs/2412.00451}{A conditional Generative Adversarial network model for the Weather4Cast 2024 Challenge}, arXiv:2412.00451. \textit{Preprint.}
  \item Presented thesis-based nowcasting work at CVIP 2025, IIT Ropar, and earlier OLR/optical-flow work at the NCVPRIPG 2023 Student Research Symposium, IIT Jodhpur.
\end{itemize}
\section{Research Software and Applied Engineering}
\project{Nilakan Programming Language}{\href{https://nilakan.jeetsh4h.workers.dev}{Documenation \& Playground}}
\begin{itemize}
  \item Implement Nilakan with Prajas Naik; language fundamentals designed by Prof. Aamod Sane.
  \item Implement the parser, static checks, interpreter, runtime diagnostics, CLI, documentation, and browser playground using TypeScript, Chevrotain, Monaco, and Playwright.
\end{itemize}
\entry{Voltek AI}{Oct 2024 -- Aug 2026}{Software Engineer (Part-time)}
\begin{itemize}
  \item Built a battery-research assistant with internal-document retrieval, web search, citations, streamed responses, and chat history across four model providers.
  \item Redesigned PostgreSQL process storage; reduced observed material-query latency from 3--20 seconds to under 500 ms with indexed tables, transactional writes, and validated backfills.
\end{itemize}
\section{Technical Skills}
\textbf{Languages:} Python, TypeScript, C++, R, SQL, Rust.\\
\textbf{ML and data:} TensorFlow/Keras, PyTorch, NumPy, OpenCV, HDF5, ConvLSTM, ConvGRU, conditional GANs, variational autoencoders, retrieval-augmented generation.\\
\textbf{Infrastructure:} Linux, GPU-server administration, air-gapped HPC, Docker, Git, GitHub Actions.
\end{document}
